\section{Smart Pointers}

\subsection{Section Overview}
Manual memory management with raw pointers (\texttt{new} and \texttt{delete}) is a major source of errors in C++, leading to memory leaks and dangling pointers. Modern C++ introduced smart pointers, which are wrapper classes that manage the lifetime of dynamically allocated memory automatically. This is a key aspect of modern C++ resource management, often referred to as RAII (Resource Acquisition Is Initialization).

\subsection{Issues with Raw Pointers}
\begin{itemize}
    \item \textbf{Forgetting to \texttt{delete}}: Leads to memory leaks.
    \item \textbf{\texttt{delete}ing too early}: Leads to dangling pointers.
    \item \textbf{\texttt{delete}ing the same memory twice}: Leads to undefined behavior.
    \item \textbf{Not using \texttt{delete[]} for arrays}: Leads to undefined behavior.
\end{itemize}
Smart pointers are designed to solve these problems.

\subsection{What is a Smart Pointer?}
A smart pointer is an object that acts like a pointer but also manages the memory it points to. When the smart pointer object goes out of scope, its destructor is automatically called, and the memory it manages is freed. This makes memory management safer and easier. The C++ Standard Library provides three types of smart pointers in the \texttt{<memory>} header: \texttt{std::unique\_ptr}, \texttt{std::shared\_ptr}, and \texttt{std::weak\_ptr}.

\subsection{Unique Pointers (unique\_ptr)}
\texttt{std::unique\_ptr} represents exclusive ownership. Only one \texttt{unique\_ptr} can point to a given resource at any time. They are very lightweight and have virtually no performance overhead compared to raw pointers.
\begin{itemize}
    \item Cannot be copied.
    \item Can be moved (ownership is transferred).
    \item The memory is automatically deallocated when the \texttt{unique\_ptr} goes out of scope.
\end{itemize}
\begin{lstlisting}
#include <memory>

int main() {
	// Use std::make_unique to create a unique_ptr (C++14+)
	std::unique_ptr<int> p1 = std::make_unique<int>(100);

	// std::unique_ptr<int> p2 = p1; // Compile Error: cannot copy

	std::unique_ptr<int> p3 = std::move(p1); // OK: transfers ownership to p3
	                                         // p1 is now nullptr

	// memory is automatically freed when p3 goes out of scope
}
\end{lstlisting}
\texttt{unique\_ptr} is the default smart pointer you should choose unless you specifically need shared ownership.

\subsection{Shared Pointers (shared\_ptr)}
\texttt{std::shared\_ptr} represents shared ownership. Multiple \texttt{shared\_ptr}s can point to the same resource. The resource is deallocated only when the last \texttt{shared\_ptr} pointing to it is destroyed. This is managed through a reference count.
\begin{itemize}
    \item Can be copied. Copying increments the reference count.
    \item The resource is deleted when the reference count becomes zero.
\end{itemize}
\begin{lstlisting}
#include <memory>

int main() {
	// Use std::make_shared to create a shared_ptr
	std::shared_ptr<int> s1 = std::make_shared<int>(100);
	{
		std::shared_ptr<int> s2 = s1; // copy, ref count is now 2
		std::cout << "Ref count: " << s1.use_count() << std::endl; // 2
	} // s2 goes out of scope, ref count is now 1
	std::cout << "Ref count: " << s1.use_count() << std::endl; // 1

} // s1 goes out of scope, ref count is 0, memory is deleted
\end{lstlisting}

\subsection{Weak Pointers (weak\_ptr)}
\texttt{std::weak\_ptr} is a non-owning smart pointer that holds a ``weak'' reference to an object managed by a \texttt{shared\_ptr}. It is used to break circular references between \texttt{shared\_ptr} objects.
\begin{itemize}
    \item Does not affect the reference count.
    \item Cannot be dereferenced directly. You must convert it to a \texttt{shared\_ptr} first.
\end{itemize}
\begin{lstlisting}
class B; // forward declaration
class A {
public:
	std::shared_ptr<B> b_ptr;
};
class B {
public:
	// std::shared_ptr<A> a_ptr; // This would create a circular reference
	std::weak_ptr<A> a_ptr; // Use weak_ptr to break the cycle
};
\end{lstlisting}

\subsection{Custom Deleters}
Smart pointers can be configured with custom deleters, which are functions or function objects that are called to release the resource. This is useful when you are managing resources other than memory allocated with \texttt{new} (e.g., file handles, network sockets).
\begin{lstlisting}
void my_deleter(int *p) {
	std::cout << "Custom deleter called" << std::endl;
	delete p;
}

int main() {
	std::shared_ptr<int> p(new int(100), my_deleter);
} // my_deleter is called when p goes out of scope
\end{lstlisting}

\subsection{Section Challenge}
Refactor the Polymorphism section challenge (the \texttt{Shape} hierarchy) to use smart pointers (\texttt{std::unique\_ptr}) instead of raw pointers. The program should have no raw pointers and no calls to \texttt{new} or \texttt{delete}.

\subsection{Section Challenge --- Solution}
\begin{lstlisting}
#include <iostream>
#include <vector>
#include <memory>

class Shape {
public:
	virtual double area() const = 0;
	virtual ~Shape() = default;
};

class Circle : public Shape {
private:
	double radius;
public:
	Circle(double r) : radius(r) {}
	double area() const override {
		return 3.14159 * radius * radius;
	}
};

class Rectangle : public Shape {
private:
	double width, height;
public:
	Rectangle(double w, double h) : width(w), height(h) {}
	double area() const override {
		return width * height;
	}
};

int main() {
	std::vector<std::unique_ptr<Shape>> shapes;

	shapes.push_back(std::make_unique<Circle>(10.0));
	shapes.push_back(std::make_unique<Rectangle>(10.0, 20.0));

	for (const auto &shape : shapes) {
		std::cout << "Area: " << shape->area() << std::endl;
	}

	// No need to delete, memory is managed automatically!

	return 0;
}
\end{lstlisting}

\subsection{Quiz 14: Section 17 Quiz}
\begin{enumerate}
    \item What is the primary problem that smart pointers solve?
    \begin{enumerate}
        \item[a)] Slow compilation times.
        \item[b)] Managing the lifetime of dynamically allocated resources to prevent leaks.
        \item[c)] Replacing all loops in a program.
        \item[d)] Storing large amounts of data.
    \end{enumerate}

    \item Which smart pointer implements exclusive ownership?
    \begin{enumerate}
        \item[a)] \texttt{std::shared\_ptr}
        \item[b)] \texttt{std::weak\_ptr}
        \item[c)] \texttt{std::auto\_ptr} (deprecated)
        \item[d)] \texttt{std::unique\_ptr}
    \end{enumerate}

    \item What is \texttt{std::weak\_ptr} primarily used for?
    \begin{enumerate}
        \item[a)] To hold the only pointer to a resource.
        \item[b)] To provide a non-owning reference to an object managed by a \texttt{shared\_ptr}, often to break circular references.
        \item[c)] As a faster version of \texttt{shared\_ptr}.
        \item[d)] To manage arrays of objects.
    \end{enumerate}
\end{enumerate}

\textbf{Answers:} 1. (b), 2. (d), 3. (b)
